% GENERATED FILE. Edit canonical lessons, then run npm run generate:books.
% canonical-source-hash: fnv1a64:4b2c65fa920bd092
% canonical-lessons: GU-C358-seva, GU-C358-tyag, GU-C358-paropkar, GU-C358-sahkar, GU-C358-nibhavvu

\chapter{Serving Others}
\label{ch:gu-serving-others}
\hlchaptermodality{358}

\begin{chapteropening}
\textbf{By the end of this chapter:} \emph{I can talk about helping others and giving things up for them.}

Complete the last lesson of chapter 358: I can talk about helping others and giving things up for them.
\end{chapteropening}

\section[\texorpdfstring{\gu{સેવા} (sevā)}{સેવા (sevā)}]{\gu{સેવા} (sevā) — service to others}

\label{lesson:GU-C358-seva}



\begin{quote}
Before the new one: say the Gujarati for a social practice, then the Gujarati for to adopt.
\end{quote}

\subsection*{You'll want to know: \gu{સેવા}}
\textbf{\gu{સેવા}} — \emph{sevā} — \textquotedblleft{}service to others\textquotedblright{}.

\begin{quote}
\gu{મીરા} \gu{દરેક} \gu{રવિવારે} \gu{દવાખાનામાં} \gu{સેવા} \gu{આપે} \gu{છે}.
\end{quote}

\begin{quote}
\emph{Mīrā darek ravivāre davākhānāmā̃ sevā āpe che.}
\end{quote}

\begin{quote}
Mira volunteers at the clinic every Sunday.
\end{quote}

A feminine noun.

\begin{cousinweb}
\textbf{\gu{સેવા}} (\emph{sevā}) is a tatsama, borrowed from Sanskrit \emph{sēvā}, \textquotedblleft{}attendance, service\textquotedblright{}; CDIAL lists an older inherited Gujarati \emph{sev}, \textquotedblleft{}attendance on\textquotedblright{}, from the same word (CDIAL 13595).
\end{cousinweb}

\begin{grammarlens}[title={What you've built}]
One more word for giving reasons and opinions.
\end{grammarlens}

\subsection*{Your turn}
\begin{itemize}
\raggedright
  \item \emph{Say it:} \emph{sevā}
  \item \emph{Say it:} \emph{sevā}, once more
  \item \emph{Say it:} \emph{Mīrā darek ravivāre davākhānāmā̃ sevā āpe che.} — Mira volunteers at the clinic every Sunday.
  \item \emph{Recall:} say the Gujarati for a social practice, then the Gujarati for to adopt, then say \emph{sevā} again
\end{itemize}

\begin{tcolorbox}[breakable,colback=teal!4,colframe=teal!35!black,title={Before you move on}]
What does \gu{સેવા} mean? (\textquotedblleft{}Service to others\textquotedblright{}.) Say it once more.
\end{tcolorbox}
\section[\texorpdfstring{\gu{ત્યાગ} (tyāg)}{ત્યાગ (tyāg)}]{\gu{ત્યાગ} (tyāg) — giving something up (for others)}

\label{lesson:GU-C358-tyag}



\begin{quote}
Before the new one: say the Gujarati for to adopt, then the Gujarati for service to others.
\end{quote}

\subsection*{You'll want to know: \gu{ત્યાગ}}
\textbf{\gu{ત્યાગ}} — \emph{tyāg} — \textquotedblleft{}giving something up (for others)\textquotedblright{}.

\begin{quote}
\gu{માતાએ} \gu{બાળકો} \gu{માટે} \gu{ઘણો} \gu{ત્યાગ} \gu{કર્યો}.
\end{quote}

\begin{quote}
\emph{mātāe bāḷko māṭe ghaṇo tyāg karyo.}
\end{quote}

\begin{quote}
The mother gave up a great deal for her children.
\end{quote}

A masculine noun.

\begin{cousinweb}
\textbf{\gu{ત્યાગ}} (\emph{tyāg}) is a tatsama, borrowed from Sanskrit \emph{tyāga}, \textquotedblleft{}leaving, abandoning, giving up\textquotedblright{} (Monier-Williams).
\end{cousinweb}

\begin{grammarlens}[title={What you've built}]
One more word for giving reasons and opinions.
\end{grammarlens}

\subsection*{Your turn}
\begin{itemize}
\raggedright
  \item \emph{Say it:} \emph{tyāg}
  \item \emph{Say it:} \emph{tyāg}, once more
  \item \emph{Say it:} \emph{mātāe bāḷko māṭe ghaṇo tyāg karyo.} — The mother gave up a great deal for her children.
  \item \emph{Recall:} say the Gujarati for to adopt, then the Gujarati for service to others, then say \emph{tyāg} again
\end{itemize}

\begin{tcolorbox}[breakable,colback=teal!4,colframe=teal!35!black,title={Before you move on}]
What does \gu{ત્યાગ} mean? (\textquotedblleft{}Giving something up (for others)\textquotedblright{}.) Say it once more.
\end{tcolorbox}
\section[\texorpdfstring{\gu{પરોપકાર} (paropkār)}{પરોપકાર (paropkār)}]{\gu{પરોપકાર} (paropkār) — doing good to others}

\label{lesson:GU-C358-paropkar}



\begin{quote}
Before the new one: say the Gujarati for service to others, then the Gujarati for giving something up.
\end{quote}

\subsection*{You'll want to know: \gu{પરોપકાર}}
\textbf{\gu{પરોપકાર}} — \emph{paropkār} — \textquotedblleft{}doing good to others\textquotedblright{}.

\begin{quote}
\gu{પરોપકાર} \gu{એ} \gu{સૌથી} \gu{મોટો} \gu{ગુણ} \gu{છે}.
\end{quote}

\begin{quote}
\emph{paropkār e sauthī moṭo guṇ che.}
\end{quote}

\begin{quote}
Doing good to others is the greatest virtue.
\end{quote}

A masculine noun.

\begin{cousinweb}
\textbf{\gu{પરોપકાર}} (\emph{paropkār}) is a tatsama, borrowed from Sanskrit \emph{paropakāra}, \textquotedblleft{}assisting others, benevolence\textquotedblright{} (Monier-Williams), built from \emph{para}, \textquotedblleft{}another\textquotedblright{}, and \emph{upakāra}, \textquotedblleft{}help\textquotedblright{}.
\end{cousinweb}

\begin{grammarlens}[title={What you've built}]
One more word for giving reasons and opinions.
\end{grammarlens}

\subsection*{Your turn}
\begin{itemize}
\raggedright
  \item \emph{Say it:} \emph{paropkār}
  \item \emph{Say it:} \emph{paropkār}, once more
  \item \emph{Say it:} \emph{paropkār e sauthī moṭo guṇ che.} — Doing good to others is the greatest virtue.
  \item \emph{Recall:} say the Gujarati for service to others, then the Gujarati for giving something up, then say \emph{paropkār} again
\end{itemize}

\begin{tcolorbox}[breakable,colback=teal!4,colframe=teal!35!black,title={Before you move on}]
What does \gu{પરોપકાર} mean? (\textquotedblleft{}Doing good to others\textquotedblright{}.) Say it once more.
\end{tcolorbox}
\section[\texorpdfstring{\gu{સહકાર} (sahkār)}{સહકાર (sahkār)}]{\gu{સહકાર} (sahkār) — cooperation}

\label{lesson:GU-C358-sahkar}



\begin{quote}
Before the new one: say the Gujarati for giving something up, then the Gujarati for doing good to others.
\end{quote}

\subsection*{You'll want to know: \gu{સહકાર}}
\textbf{\gu{સહકાર}} — \emph{sahkār} — \textquotedblleft{}cooperation\textquotedblright{}.

\begin{quote}
\gu{બધાના} \gu{સહકારથી} \gu{અમે} \gu{કામ} \gu{સમયસર} \gu{પતાવ્યું}.
\end{quote}

\begin{quote}
\emph{badhānā sahkārthī ame kām samaysar patāvyũ.}
\end{quote}

\begin{quote}
With everyone's cooperation we finished the work on time.
\end{quote}

A masculine noun.

\begin{cousinweb}
\textbf{\gu{સહકાર}} (\emph{sahkār}) is a tatsama, borrowed from Sanskrit \emph{sahakāra}, \textquotedblleft{}acting with, cooperation\textquotedblright{} (Monier-Williams).
\end{cousinweb}

\begin{grammarlens}[title={What you've built}]
One more word for giving reasons and opinions.
\end{grammarlens}

\subsection*{Your turn}
\begin{itemize}
\raggedright
  \item \emph{Say it:} \emph{sahkār}
  \item \emph{Say it:} \emph{sahkār}, once more
  \item \emph{Say it:} \emph{badhānā sahkārthī ame kām samaysar patāvyũ.} — With everyone's cooperation we finished the work on time.
  \item \emph{Recall:} say the Gujarati for giving something up, then the Gujarati for doing good to others, then say \emph{sahkār} again
\end{itemize}

\begin{tcolorbox}[breakable,colback=teal!4,colframe=teal!35!black,title={Before you move on}]
What does \gu{સહકાર} mean? (\textquotedblleft{}Cooperation\textquotedblright{}.) Say it once more.
\end{tcolorbox}
\section[\texorpdfstring{\gu{નિભાવવું} (nibhāvvũ)}{નિભાવવું (nibhāvvũ)}]{\gu{નિભાવવું} (nibhāvvũ) — to keep up (a promise or a duty)}

\label{lesson:GU-C358-nibhavvu}



\begin{quote}
Before the new one: say the Gujarati for doing good to others, then the Gujarati for cooperation.
\end{quote}

\subsection*{You'll want to know: \gu{નિભાવવું}}
\textbf{\gu{નિભાવવું}} — \emph{nibhāvvũ} — \textquotedblleft{}to keep up (a promise or a duty)\textquotedblright{}.

\begin{quote}
\gu{અરુણ} \gu{હંમેશા} \gu{પોતાનું} \gu{વચન} \gu{નિભાવે} \gu{છે}.
\end{quote}

\begin{quote}
\emph{Aruṇ hammeśā potānũ vacan nibhāve che.}
\end{quote}

\begin{quote}
Arun always keeps his word.
\end{quote}

Past, with \gu{તેણે} (he or she …): \gu{તેણે} \gu{પોતાની} \gu{ફરજ} \gu{નિભાવી}, \emph{teṇe potānī pharaj nibhāvī}.

\begin{cousinweb}
\textbf{\gu{નિભાવવું}} (\emph{nibhāvvũ}) is inherited from Sanskrit \emph{nírvahati}, \textquotedblleft{}leads out of, accomplishes\textquotedblright{}; CDIAL lists the Gujarati verb (CDIAL 7397).
\end{cousinweb}

\begin{grammarlens}[title={What you've built}]
One more word for giving reasons and opinions.
\end{grammarlens}

\subsection*{Your turn}
\begin{itemize}
\raggedright
  \item \emph{Say it:} \emph{nibhāvvũ}
  \item \emph{Say it:} \emph{nibhāvvũ}, once more
  \item \emph{Say it:} \emph{Aruṇ hammeśā potānũ vacan nibhāve che.} — Arun always keeps his word.
  \item \emph{Recall:} say the Gujarati for doing good to others, then the Gujarati for cooperation, then say \emph{nibhāvvũ} again
\end{itemize}

\begin{tcolorbox}[breakable,colback=teal!4,colframe=teal!35!black,title={Before you move on}]
What does \gu{નિભાવવું} mean? (\textquotedblleft{}To keep up (a promise or a duty)\textquotedblright{}.) Say it once more.
\end{tcolorbox}
